\subsection*{2. Acoplamientos definidos por una biextensión}

2.1. Sean $P$ y $Q$ dos grupos abelianos de $T$, y sea $n>0$ un
entero. Recordemos que $\,{}_nP$ denota el núcleo de la multiplicación
por $n\cdot\id_P$. Definiremos un homomorfismo canónico

%%%%%%%%%%  scan idx 238  |  folio 227  %%%%%%%%%%

\begin{equation}\tag{2.1.1}
 \alpha=\alpha^n_{P,Q}:
 {}_nP\otimes{}_nQ\longrightarrow\uTor_1(P,Q),
\end{equation}
que será el compuesto
\begin{equation}\tag{2.1.2}
 {}_nP\otimes{}_nQ
 \longrightarrow \uTor_1({}_nP,{}_nQ)
 \longrightarrow \uTor_1(P,Q),
\end{equation}
donde la primera flecha es $\alpha^n_{{}_nP,{}_nQ}$ y la segunda está
inducida por las inclusiones $\,{}_nP\longrightarrow P$ y
$\,{}_nQ\longrightarrow Q$. Quedamos así reducidos a definir (2.1.1)
cuando $nP=0$ y $nQ=0$, es decir, cuando $P$ y $Q$ son
$\mathbb Z/n\mathbb Z$-módulos.

Supongamos solamente que $nQ=0$. Tomando una resolución plana
$L_\bullet$ de $P$ (como $\mathbb Z$-módulo), se obtienen los
isomorfismos canónicos
\[
 P\mathbin{\overset{\mathbb L}{\otimes}_{\mathbb Z}}Q
 \simeq L_\bullet\otimes_{\mathbb Z}Q
 \simeq
 (L_\bullet\otimes_{\mathbb Z}\mathbb Z/n\mathbb Z)
 \otimes_{\mathbb Z/n\mathbb Z}Q,
\]
que pueden escribirse como un isomorfismo en la categoría derivada
\begin{equation}\tag{2.1.3}
 P\mathbin{\overset{\mathbb L}{\otimes}_{\mathbb Z}}Q
 \simeq
 \left(P\mathbin{\overset{\mathbb L}{\otimes}_{\mathbb Z}}
       \mathbb Z/n\mathbb Z\right)
 \mathbin{\overset{\mathbb L}{\otimes}_{\mathbb Z/n\mathbb Z}}Q.
\end{equation}
Esto da lugar a la sucesión espectral homológica de tipo Künneth
(cf. SGA 1 IV 5.2)
\begin{equation}\tag{2.1.4}
 \uTor^{\mathbb Z}_*(P,Q)
 \Longleftarrow
 E^2_{p,q}
 =
 \uTor^{\mathbb Z/n\mathbb Z}_p
 \bigl(\uTor^{\mathbb Z}_q(P,\mathbb Z/n\mathbb Z),Q\bigr).
\end{equation}
Recordemos también los isomorfismos canónicos
\[
 \uTor^{\mathbb Z}_0(P,\mathbb Z/n\mathbb Z)
 =P_n\mathrel{\mathop:}=P/nP,
 \qquad
 \uTor^{\mathbb Z}_1(P,\mathbb Z/n\mathbb Z)\simeq{}_nP,
\]
y
\[
 \uTor^{\mathbb Z}_i(P,\mathbb Z/n\mathbb Z)=0
 \qquad\text{para }i\geq2.
\]
Se deducen calculando estos objetos mediante la resolución plana
canónica de $(\mathbb Z/n\mathbb Z)_T$:
\[
\begin{tikzcd}[column sep=large,row sep=large]
 0\arrow[r]&
 \mathbb Z_T\arrow[r,"n"]\arrow[d,equal]&
 \mathbb Z_T\arrow[r]\arrow[d,equal]&
 0\\
 &L_1&L_0&
\end{tikzcd}
\]

%%%%%%%%%%  scan idx 239  |  folio 228  %%%%%%%%%%

cuya aumentación
$L_0\longrightarrow(\mathbb Z/n\mathbb Z)_T$ envía $1$ en $1$.
Por consiguiente, la sucesión espectral (2.1.4) se reduce
esencialmente a la sucesión exacta
\begin{equation}\tag{2.1.5}
\begin{split}
0\longrightarrow{}&
 \uTor^{\mathbb Z/n\mathbb Z}_2(P_n,Q)
 \longrightarrow{}_nP\otimes Q
 \xrightarrow{\alpha}\uTor^{\mathbb Z}_1(P,Q)\\
\longrightarrow{}&
 \uTor^{\mathbb Z/n\mathbb Z}_1(P_n,Q)
 \longrightarrow0,
\end{split}
\end{equation}
donde $P_n=P\otimes_{\mathbb Z}\mathbb Z/n\mathbb Z=P/nP$.
El homomorfismo buscado $\alpha$ es el homomorfismo de borde
$E^2_{0,1}\longrightarrow\uTor^{\mathbb Z}_1$ de esta sucesión
exacta. Como la sucesión es manifiestamente funtorial en $P$ y $Q$,
la funtorialidad en $P$ muestra que el homomorfismo considerado es en
efecto el compuesto
\[
 {}_nP\otimes Q
 \longrightarrow\uTor_1({}_nP,Q)
 \longrightarrow\uTor_1(P,Q),
\]
donde el primer homomorfismo es el $\alpha$ definido anteriormente,
con $(P,Q)$ reemplazado por $({}_nP,Q)$, y el segundo está inducido
por la inclusión $\,{}_nP\longrightarrow\sisi{Q}{P}$.

2.1.6. Esto define, para $P$ y $Q$ cualesquiera, un homomorfismo
(2.1.1) que satisface la condición de factorización (2.1.2). Además,
cuando $nQ=0$, la sucesión exacta (2.1.5) muestra que $\alpha$ es un
isomorfismo siempre que $P_n$ sea plano sobre
$\mathbb Z/n\mathbb Z$. Esto ocurre en particular si $P_n=0$, es
decir, si $nP=P$ y $P$ es $n$-divisible, o si el propio $P$ es un
$\mathbb Z/n\mathbb Z$-módulo plano. También ocurre si $Q$ es un
$\mathbb Z/n\mathbb Z$-módulo plano.

2.1.7. Explicitemos el cálculo del homomorfismo $\alpha$ de (2.1.5).
Para determinarlo basta determinar, para toda sección $q$ de $Q$
(sobre un objeto $S$ de $T$; tomemos el objeto final para simplificar),
el homomorfismo correspondiente
\[
 {}_nP\longrightarrow\uTor^{\mathbb Z}_1(P,Q).
\]
La sección $q$ puede verse como la imagen de la sección $1$ de
$(\mathbb Z/n\mathbb Z)_T$ por el homomorfismo unívocamente
determinado
\[
 u_q:(\mathbb Z/n\mathbb Z)_T\longrightarrow Q.
\]
La funtorialidad de $\alpha$ en $Q$ nos reduce, por tanto, al caso
$Q=(\mathbb Z/n\mathbb Z)_T$, con $q$ su sección unidad.

%%%%%%%%%%  scan idx 240  |  folio 229  %%%%%%%%%%

En este caso, identifiquemos
\[
 \uTor^{\mathbb Z}_1(P,Q)
 =
 \uTor^{\mathbb Z}_1(P,\mathbb Z/n\mathbb Z)
\]
con $\,{}_nP$, como se explicó después de (2.1.4). El homomorfismo
buscado es entonces la identidad. Equivalentemente, el homomorfismo
de borde
\[
 E^2_{0,q}
 =
 \uTor^{\mathbb Z}_q(P,\mathbb Z/n\mathbb Z)
 \otimes_{\mathbb Z/n\mathbb Z}\mathbb Z/n\mathbb Z
 \longrightarrow
 \uTor^{\mathbb Z}_q(P,\mathbb Z/n\mathbb Z)
\]
es la identidad. Esto se deduce inmediatamente de las definiciones.
La afirmación análoga sigue siendo válida al reemplazar $\mathbb Z$ y
su cociente $\mathbb Z/n\mathbb Z$ por cualquier anillo $A$ sobre
$T$ y cualquier anillo cociente $A/I$.

2.1.8. El homomorfismo $\alpha$ de (2.1.1) también puede determinarse
describiendo, para cada par de secciones $p$ de $\,{}_nP$ y $q$ de
$\,{}_nQ$ (sobre un objeto $S$ de $T$, que tomamos de nuevo final para
simplificar la notación), la sección de $\uTor^{\mathbb Z}_1(P,Q)$
determinada por $(p,q)$. Por funtorialidad en $P$ y $Q$, como antes,
quedamos reducidos al caso
\[
 P=Q=(\mathbb Z/n\mathbb Z)_T,
\]
con $p$ y $q$ iguales a la sección unidad. Basta, pues, determinar
\[
 \alpha:
 (\mathbb Z/n\mathbb Z)_T\otimes(\mathbb Z/n\mathbb Z)_T
 \longrightarrow
 \sisi{\uTor^{\mathbb Z}_1
 (\mathbb Z/n\mathbb Z)_T,(\mathbb Z/n\mathbb Z)_T)}%
 {\uTor^{\mathbb Z}_1
 \bigl((\mathbb Z/n\mathbb Z)_T,(\mathbb Z/n\mathbb Z)_T\bigr)}.
\]
La descripción de 2.1.7 muestra que este es el isomorfismo
\begin{equation}\tag{2.1.8.1}
 (\mathbb Z/n\mathbb Z)_T
 \xrightarrow{\sim}
 \sisi{\uTor^{\mathbb Z}_1
 (\mathbb Z/n\mathbb Z)_T,(\mathbb Z/n\mathbb Z)_T)}%
 {\uTor^{\mathbb Z}_1
 \bigl((\mathbb Z/n\mathbb Z)_T,(\mathbb Z/n\mathbb Z)_T\bigr)}
\end{equation}
deducido de la resolución plana canónica
\begin{equation}\tag{2.1.8.2}
 0\longrightarrow\mathbb Z_T
 \xrightarrow{n}\mathbb Z_T
 \longrightarrow0
\end{equation}

%%%%%%%%%%  scan idx 241  |  folio 230  %%%%%%%%%%

del \underline{primer} factor
$P=(\mathbb Z/n\mathbb Z)_T$ que aparece en $\uTor_1$.
